\section{Phase 6: two-dimensional regime map}
\label{sec:phase6}

\subsection{Experiment 6A: scheduling--heterogeneity sweep}
\paragraph{Objective.}
Determine when each architecture is warranted rather than claiming that M4 is
universally best.

\paragraph{Setup.}
Scale structural separation using
\begin{equation}
 \vartheta_c(\gamma_h)=\vartheta_0+\gamma_h(\vartheta_c-\vartheta_0),
 \qquad \gamma_h\in\{0,0.5,1,1.5\},
\end{equation}
and use speed envelopes
$[18,22]$, $[15,25]$, and $[10,30]\,\mathrm{m\,s^{-1}}$. Compare M1--M4 in every
cell using model error, closed-loop performance, worst-family result, and
complexity.

\paragraph{Predeclared expectation.}
Global LTI is sufficient for low heterogeneity and narrow scheduling variation.
Global LPV dominates for broad scheduling variation but low heterogeneity.
Clustered LTI dominates for high heterogeneity over a narrow speed range. Clustered
LPV becomes necessary when both dimensions are large.

\paragraph{Planned evidence and decision gate.}
The main figure is a two-dimensional map of the best architecture and the relative
advantage of M4. The map must show coherent transitions rather than one method
winning due to numerical noise. Unexpected regions are investigated, not smoothed
away.

\paragraph{Results, interpretation, and decision.}
Pending.
